\documentclass[11pt,a4paper]{article}
\usepackage[margin=1in]{geometry}
\usepackage{booktabs}
\usepackage{amsmath}
\usepackage[hidelinks]{hyperref}
\usepackage{parskip}
\usepackage{lmodern}
\usepackage[T1]{fontenc}

\title{\emph{Cis}-pQTL Mendelian Randomization and Phase~III Drug-Target Outcomes:\\
Corrected Analysis Report}
\author{Elliot Tower\\ \texttt{elliot@elliottower.ai}}
\date{17 August 2026}

\begin{document}
\maketitle

\begin{abstract}
\noindent
This report accompanies the deposited analysis archive. It states the corrected results for
a pre-registered evaluation of \emph{cis}-pQTL Mendelian randomization as a predictor of
Phase~III trial success, records the outcome-GWAS corrections applied to the underlying
data, and reports the test of the registered mechanism hypothesis. The classifier
discriminates above chance overall, with a positive predictive value of 0.588 against a
base rate of 0.308. The registered prediction that performance differs by drug mechanism is
not supported at this sample size, and the design cannot resolve it at any sample size the
instrument universe permits. No manuscript accompanies this deposit.
\end{abstract}

\section{Scope}

The evaluation covers 143 drug-target-indication pairs across 91 genes and 24 diseases.
Trial outcomes were adjudicated from FDA and EMA approval records and ClinicalTrials.gov
before any Mendelian randomization estimate was computed, under the staged architecture
recorded in \texttt{protocol/DEVIATION\_LOG.md}. The candidate set was frozen and hashed at
that point; \texttt{data/frozen\_candidates\_v34.csv} hashes to \texttt{c428c733\ldots},
matching the value recorded in the deviation log.

The classifier is a single-SNP Wald ratio on a sentinel \emph{cis}-pQTL, with $p < 0.05$
predicting SUCCESS. It has no tunable parameters and was fixed in protocol V3.1.

\section{Results}

\begin{table}[h]
\centering
\begin{tabular}{lrrrl}
\toprule
Stratum & $n$ & Balanced accuracy & 90\% CI & Criterion \\
\midrule
Primary & 143 & 0.5783 & [0.5223, 0.6364] & met \\
Effective $N \geq 2{,}000$ & 134 & 0.5810 & [0.5243, 0.6415] & met \\
Abundance-modulating & 65 & 0.6165 & [0.5309, 0.7040] & met \\
Activity-blocking & 75 & 0.5136 & [0.4554, 0.5826] & not met \\
\bottomrule
\end{tabular}
\caption{Balanced accuracy by stratum, with 90\% bootstrap intervals over 10{,}000
resamples. The pre-registered falsification criterion is a lower bound above 0.50.
Base rate 0.308 SUCCESS.}
\end{table}

\vspace{-0.5em}
Sensitivity is 0.227 and specificity 0.929. The classifier fires on 17 of 143 pairs and is
correct on 10 of those, giving a positive predictive value of 0.588 against a base rate of
0.308, a 1.9-fold lift. That figure rests on 17 positive calls drawn from a small number of
genes and is reported with that caveat throughout. It is consistent with the roughly
two-fold lift for genetically supported targets reported by Nelson et al. (2015) and King
et al. (2019), which are large replicated estimates; the agreement here is concordance with
an established benchmark rather than independent corroboration of it.

The negative predictive value is 0.730 against a failure base rate of 0.692, a lift of
1.05. A positive result carries information; a null result carries almost none.

\section{The mechanism hypothesis}

Protocol V3.4 registered the prediction that the classifier performs better for
abundance-modulating targets, where the drug acts on the circulating protein the assay
measures, than for activity-blocking targets, where it does not.

The observed difference is $+0.1028$. Gene-clustered permutation of the mechanism label
over 20{,}000 iterations gives $p = 0.2144$; a gene-clustered bootstrap on the difference
gives a 90\% interval of $[-0.0198, +0.2262]$. Genes are the resampling unit because pairs
sharing a gene share an exposure GWAS and an instrument. \textbf{The prediction is not
supported at this sample size.}

The difference is also sensitive to specification. It falls to $+0.0400$ among the 100
pairs whose outcome GWAS can be identified, to $+0.0148$ outside immune-mediated disease,
and reverses to $-0.0349$ within UKB-PPP instruments. Two further stratifications registered
in the same protocol produce differences of comparable size and comparable
non-significance: disease area $+0.0941$ ($p = 0.281$) and instrument source $+0.0757$
($p = 0.385$).

Detecting a difference of the observed size at 80\% power requires approximately 345 to 435
analyzed pairs. The applicability funnel places the entire universe of Phase~III
drug-target pairs carrying a plasma \emph{cis}-pQTL instrument at 211. The hypothesis is
not resolvable with data that currently exist.

\section{Outcome-GWAS corrections}

An audit of every accession found eight errors. Five pointed at an unrelated trait: chronic
kidney disease to a bladder cancer GWAS, systemic lupus erythematosus to serum copper,
autism spectrum disorder to oral and pharyngeal cancer, melanoma to waist circumference,
and hypercholesterolemia to an LDL cholesterol GWAS against an LDL-lowering endpoint. Three
had been withdrawn from the OpenGWAS index after they were queried.

A separate \texttt{DISEASE\_GWAS} mapping in \texttt{code/expand\_v5.py} supplied two
further diseases and had never been audited. Juvenile idiopathic arthritis used a 216-case
substitute where the deviation log pre-specified a 2{,}816-case source, and the
substitute's controls include UK Biobank participants while five of the six pairs use
UKB-PPP instruments, contrary to the overlap policy in force.

Every correction restores a source named in the deviation log, with one exception: no
glioma GWAS matching the pre-specified consortium exists in OpenGWAS, and the pre-declared
FinnGen fallback was used. The two mappings are now a single shared constant.

\textbf{The corrections did not change the conclusion.} The mechanism difference was
$+0.0955$ before correction, $+0.1028$ after, and $+0.1029$ after LD-proxy recovery was
applied symmetrically to both the original and the corrected accessions. What changed the
conclusion was testing the difference rather than inferring it from two intervals.

\section{Limits on provenance}

Forty of the 143 pairs take their estimate from the EpiGraphDB MR-EvE catalog. The catalog
records source, method, SNP count and rsID, and does not record which outcome GWAS was
used. For those pairs the outcome phenotype, its case and control counts, and any sample
overlap with the exposure cohort cannot be verified. They reproduce exactly from the
deposited catalog file, which is included here and had not previously been public.

After LD-proxy recovery, thirteen pairs remain unrecovered because their instrument is
absent from the corrected outcome GWAS and has no proxy at $r^2 \geq 0.6$. Those
instruments have a median minor allele frequency of 0.100 against 0.215 for retained pairs,
on nine and eighty-nine observations respectively. The analyzed set is depleted of
low-frequency instruments in pancreatic cancer, chronic kidney disease and juvenile
idiopathic arthritis.

The sample-overlap sensitivity analysis is not evaluable. One flagged pair survives the
corrections, and a comparison against a single-pair stratum estimates nothing.

\section{Archive contents}

The deposit holds the pre-registration chain V1 through V7 with amendments and freeze
records, the corrected analysis and its caches, the prior uncorrected analyzed set for
comparison, the frozen candidate and adjudication data, the EpiGraphDB catalog, the
analysis code and the correction harness, and a SHA-256 manifest.

Protocols V1 and V3 do not verify against their recorded hashes. Both hash files name
pre-rename filenames and both documents entered the repository already renamed, so those
records describe files that never existed under those names. Both are superseded.

No manuscript is deposited. The current draft asserts a mechanism-dependent difference that
these data do not support.

\end{document}
